\cotefile{134-2}
\scanpage{83}
\mspage{\unsure{I|1}}
% archive sheet number: 82
\scanpage{84}
\begin{center}
APPENDIX \deleted{\unsure{to}}TO Chapter I : three letters to Larry Breen
\end{center}

\deleted{\unsure{I}} In this appendix, I am encluding three letters to Larry Breen, dated 5.2 , 17.2 and 17-19.7. 1975. These letters are written in French, and the first two are reproduced textually, whereas the third appears here in English translation. I am grateful to Ronnie Brown who took the trouble last year to make such a translation from a hardly legible \deleted{handwri} copy of the handwritten letter to Larry Breen, with the assistance of Larry Breen himself and J.L. Loday. I am now using his translation rather than the original letter, which no printer could possibly decipher ! Also, for the second letter I am using a typed copy made in 1975 by Larry Breen (who presumably had difficulties too deciphering the handwriting), Thanks are due to him for his interest and patience (with someone like me, very unknowledgeable in standard homotopy techniques), \deleted{\unsure{as it}}\inserted{\handwritten{which}} appeared \deleted{from}\inserted{\handwritten{in}} the letters I got from him in response and his verbal explanations on related matters, \inserted{\handwritten{as well as in}} the trouble he took in \deleted{copying} retyping \inserted{\handwritten{that letter}} and sending me copies of my own letters to him, and allowing me to reproduce them \deleted{\unsure{s}}in th\deleted{\unsure{a}}is volume of P\deleted{u}\handwritten{u}rsuing Stacks. My only present contribution to this set of letters is adding a few comments (in the notes (\handwritten{22}) to (\handwritten{31})), and correcting some inaccuracies in the English translation (due mainly to my handwriting ...) Also, I skipped the beginning of the first letter, which does'nt seem to me of general relevance.

The first two letters are an attempt to explain to Larry Breen (who has a wide background in algebraic geometry and homological and homotopica\unsure{l} algebra) some of the main points of the programme I had in mind\deleted{\unsure{,|x}} around the notions of $n$-categories and $n$-stacks\deleted{,} \handwritten{(}which is what I am supposed to be pursuing now in my present work ``Pursuing Stacks''\handwritten{)}. They were written under the impetus of the \inserted{new} intuition \deleted{which} (new to me at any rate), which then had just appeared to me, namely that (non strict) $n$-groupoids should model (in a suitable sense) $n$-truncated homotopy types.\handwritten{\(\mathord{/\mkern-5mu/}\)} The third letter, \inserted{\handwritten{written} in} answer to a number of questions in Larry Breen's response to the first two, is of a wider scope. A large part of the letter outlines (very ske\deleted{\unsure{ch}}\handwritten{\unsure{tch}}yly) some main points of a duality program (including a cohomological formulation of ``geometric'' local and global classfield theory), which appears here for the first time in print, \deleted{and} going back to the end of the fifties. It gives also some hints about the need of a framework of ``tame topology'' \deleted{\unsure{and}}\unsure{sui}table for writing up a ``dév\deleted{\unsure{is}}\unsure{ii|is}sage theory'' of stratified spaces, and for \deleted{\unsure{defining}}\inserted{\unsure{r|w}orking} with étale tubular neighbourhoods, for the common purpos\unsure{e} of coming to grips \deleted{of}\inserted{\handwritten{with}} a suitable notion of ``fine homotopy type'' of a ``tame\unsure{''} topological space or \inserted{\handwritten{of}} a scheme, in terms of the \deleted{ordered set}\inserted{\handwritten{inverse system}} of ``\deleted{\unsure{h}}indexed homotopy types'' corresponding to equisingular stratifications .
% archive sheet number: 83
\scanpage{85}
% A single diagonal cancellation runs from above the heading to below item 24.
% Each line within its extent is set as deleted, including the locally retyped lines.
\noindent\deleted{RECOLTES ET SEMAILLES \deleted{\unsure{-}\hspace{1em}S}}\par
\noindent\deleted{SOMMAIRE}\par
\noindent\deleted{I \quad Travail et découverte}\par
\noindent\hspace*{1em}\deleted{1. L'enfant et le Bon Dieu}\par
\noindent\hspace*{1em}\deleted{2. Erreur et découverte}\par
\noindent\hspace*{1em}\deleted{3. Les inavouables labeurs}\par
\noindent\hspace*{1em}\deleted{4. Infaillibilité (des autres) et mépris (de soi)}\par
\noindent\deleted{II \quad Le rêve et le Rêveur}\margin{\handwritten{\tikz[baseline]{\draw (0,0)--(0,.32)--(.8,.32);}}}\par
\noindent\hspace*{1em}\deleted{5. Le rêve interdit}\par
\noindent\hspace*{1em}\deleted{6. Le Rêveur}\par
\noindent\hspace*{1em}\deleted{7. L'héritage de Galois}\par
\noindent\hspace*{1em}\deleted{8. Rêve et démonstration}\par
\noindent\deleted{III \quad \deleted{\unsure{n}}\handwritten{N}aissance de la \unsure{C|c}rainte}\par
\noindent\hspace*{1em}\deleted{9. L'étranger bienvenu}\margin{\handwritten{\tikz[baseline]{\draw (.32,.85)--(.32,.73)--(.08,.73) .. controls (0,.73) and (0,.69) .. (0,.57) --(.06,.14) .. controls (.07,.08) and (.17,.12) .. (.17,.02) --(.18,-.17);}}}\par
\noindent\hspace*{1em}\deleted{10. La ``Communauté mathématique'' : fiction et réalité\deleted{\unsure{,|.}}}\par
\noindent\hspace*{1em}\deleted{11. Rencontre avec Claude Chevalley\unsure{,}- ou: liberté et bons sentiments}\par
\noindent\hspace*{1em}\deleted{12. Le mérite et le mépris}\par
\noindent\hspace*{1em}\deleted{13. Force et épaisseur}\par
\noindent\hspace*{1em}\deleted{14. Naissance de la crainte}\par
\noindent\hspace*{1em}\deleted{15. Récoltes et semailles}\par
\noindent\deleted{\deleted{IV Le double Visage}}\par
\noindent\deleted{IV \quad Le double \deleted{\unsure{v}}Visage}\par
\noindent\hspace*{1em}\deleted{16. Marais et premiers rangs}\par
\noindent\hspace*{1em}\deleted{17. Terry Mirkil}\par
\noindent\hspace*{1em}\deleted{18. Vingt ans de fatuité, ou: l'ami infatigable}\par
\noindent\hspace*{1em}\deleted{19. Le monde sans amour}\par
\noindent\hspace*{1em}\deleted{20. Un monde sans conflit ?}\par
\noindent\hspace*{1em}\deleted{21. Un secret de Polichinelle bien gardé}\par
\noindent\hspace*{1em}\deleted{22. Bourbaki,\deleted{\unsure{-}} ou ma grande chance - et son revers}\par
\noindent\hspace*{1em}\deleted{23. De Profundis}\par
\noindent\hspace*{1em}\deleted{24. Mes adieux,- ou: les étrangers}\par
% archive sheet number: 84
\scanpage{86}
\mspage{\handwritten{63}}

(\textsuperscript{12}!) This is nonsense, as one sees already in the following picture where $f\ f''=\mathrm{S}_{\unsure{o|0}}$ :

(\textsuperscript{13}!) This is nonsense again for $n\quad 1$, see PS at the end of section 13.

(\textsuperscript{14}!) Same mistake as \deleted{in} the one noticed in the previous note. The fibered products exist in $\underline{\mathrm{C}}_{o}$ only, and \uline{these} should be preserved by the functors under consideration. Thus the ``universal problem'' \deleted{\unsure{(in}}\unsure{'}has to be rephrased somewhat ...

(\textsuperscript{15}) I guess they are\unsure{.} not equivalent, even when restricting to objects which are groupoids (i.e. so-called Gr-categories).

(\textsuperscript{16}) An object of $\mathrm{M}$ is called fibrant (resp. trivial) it the map from it to the final object is fibrant (resp. a weak equivalence).

(\textsuperscript{17}) This seems doubtful, unless all objects of $\mathrm{M}$ are fibrant (which is true for the most familiar cases I was having in mind). But even assuming this, we still need a unicity statement (in a suitable sense) for the functors $\underline{\mathrm{B}}_{o}\quad \mathrm{M}$ thus obtained, in order to be sure that the correspon\deleted{d}ding functors $\mathrm{M}\quad(\mathrm{Ho})$ are all canonically isomorphic. These kind of questions may be viewed as closely related to the question of existence (and unicity) of ``test functors'' for a given test category $\mathrm{A}$ (here, the category of ``standard hemispheres'' ) into a given asphericity or contractibility structure, as discussed in section 90 (without yet getting there \deleted{\unsure{the}}\inserted{any} clear-cut handy existence and unicicity theorems which are to be \deleted{\unsure{hped}} hoped for). The set-up of asphericity and contractibility structures, which has been worked out in the months following, has been gradually replacing Quillen's approach to homotopy models. (I'll have to come back in due course to theq\deleted{\unsure{z}}uestion of the relationship between the two approaches, which deserves to be understood.)

(\textsuperscript{18}) The word ``gerb'' here stands \deleted{for} as a ``translation'' of the F\shortstack{\scriptsize\unsure{''}\\r\\[-.4ex]\scriptsize\unsure{\textasciitilde}}ench word ``gerbe'', as used in Giraud's book on non commutative cohomological algebra. With the terminology of next section, we would call it rather a ``stack of groupoids'' or ``Gr-stack'' (with a specification by $n$ or $\infty$ if needed !).

(\textsuperscript{19}\deleted{\unsure{)}}\unsure{!}) This is false, as seen below.

(\textsuperscript{20}\deleted{)}!) This descr\deleted{\unsure{p}}iption is dubious, as it may give categories $\underline{\mathrm{C}}_{n}$ which are larger than the ones we want.

(\textsuperscript{21}!) It seems dubious however that the mere category structure of $\underline{\mathrm{C}}$ will allow us to recover the \deleted{\unsure{su}} ``primitive'' subcategory $\underline{\mathrm{C}}_{o}$ , and it looks safer to add the latter as an extra structure to $\underline{\mathrm{C}}$ .
% archive sheet number: 85
\scanpage{87}
% One long pen stroke through the leaf crosses the index and its explanatory note.
\noindent\deleted{NOTES (\textsuperscript{*})}

\diagram{\begin{tabular}{@{}r@{\enspace}p{.61\linewidth}@{\enspace}r@{\enspace}l@{}}
\deleted{1.} & \deleted{Mes amis de Survivre et Vivre} & \deleted{11} & \deleted{\handwritten{(\unsure{5|0})}} \\
\deleted{2.} & \deleted{Aldo Andreotti, Ionel Bucur} & \deleted{14} & \deleted{\handwritten{(11)}} \\
\deleted{3.} & \deleted{Jésu\deleted{\unsure{s}}s et les douze apôtres} & \deleted{25} & \deleted{\handwritten{(\unsure{19|18})}} \\
\deleted{4.} & \deleted{L'Enfant et le \deleted{m}\handwritten{M}aître} & \deleted{26} & \deleted{\handwritten{(23)}} \\
\deleted{5.} & \deleted{La peur de jouer} & \deleted{29} & \deleted{\handwritten{(23\textsuperscript{\unsure{ii|''}})}} \\
\deleted{6.} & \deleted{Les deux frères} & \deleted{29} & \deleted{\handwritten{(23\textsuperscript{\unsure{iii|'''}})}} \\
\deleted{7.} & \deleted{Causes d'un échec} & \deleted{31} & \deleted{\handwritten{(23\textsuperscript{iv})}} \\
\deleted{8.} & \deleted{Consensus déontologique - et contrôle de l'information} & \deleted{32} & \deleted{\handwritten{(25)}} \\
\deleted{9.} & \deleted{Le ``snobisme des jeunes'', ou les défenseurs de la pureté} & \deleted{33} & \deleted{\handwritten{(27)}} \\
\deleted{10.} & \deleted{Cent fers dans le feu, ou: rien ne sert de sécher !} & \deleted{36} & \deleted{\handwritten{(32)}} \\
\deleted{11.} & \deleted{L'étreinte impuissante.} & \deleted{37} & \deleted{\handwritten{(34)}} \\
\deleted{12.} & \deleted{La visi\deleted{\unsure{t|i}}te, ou le rêve malvenu} & \deleted{45} & \deleted{\handwritten{(40)}} \\
\deleted{13.} & \deleted{Krishna\deleted{\unsure{n}}murti, ou \deleted{\unsure{une}}\inserted{\handwritten{la}} libération devenue entrave} & \deleted{45} & \deleted{\handwritten{(41)}} \\
\deleted{14.} & \deleted{L'arrachement salutaire} & \deleted{45} & \deleted{\handwritten{(42)}} \\
\deleted{15.} & \deleted{Mes orphelins} & \deleted{50} & \deleted{\handwritten{(45\deleted{\unsure{g}}\unsure{a})}} \\
 & \deleted{Refus d'un héritage, ou le prix d'une contradiction\unsure{,}} & \deleted{50} & \deleted{\handwritten{(45b)}} \\
\deleted{16.} & \deleted{On n'arrête pas le progrès !} & \deleted{50} & \deleted{\handwritten{(47)}} \\
\end{tabular}}

\deleted{(\textsuperscript{*}) Voici une liste des notes les plus longues ou les plus substantielles, par ordre d'appa\unsure{ri}tion dans le texte. Les numéros placés après le \deleted{\unsure{no|nom}} nom de chaque note (lequel, je l'espère, \deleted{\unsure{les|le}}\unsure{la|les|le} rend aisément reconnaissable) indiquent celui de la section (entre 1 et 50) où la note figure, suivi de celui de la note elle-même. Une des notes, plus longue que la longue section à laquelle elle se rapporte, a été subdivisée après coup en deux parties, auxquelles il m'a paru utile de donner des noms séparés.}
% archive sheet number: 86
\scanpage{88}
\mspage{\handwritten{64}}

(\textsuperscript{22}) For the notion of a ``lien'' or ``tie'', \deleted{\unsure{see}}which is one of the main ingredients of the non commutative cohomology panoply of Giraud's theory, I refer to his book (Springer Grundlehren 179, 1971). \inserted{A} Picard categor\deleted{ies}y \deleted{\unsure{are categories}}\inserted{is a groupoi\unsure{d}} endowed with an operation $\unsure{\otimes|\oplus}$ \deleted{which makes it \unsure{ressemble|resemble} to an} together with associativity, unity and commutativity data for this operation, which make it reseumble to a commutative group. \inserted{A} ``Champ\deleted{\unsure{s}} de Picard'' or ``Picard stacks'' \deleted{are}\inserted{is} defined accordingly, by relativizing over an arbitrary space or topos (replacing the groupoid by a stack of groupoids over this topos). The necessary generally nonsense on these is devel\unsure{p|o}ped rathe\unsure{r} carefully in an exposé of Deligne in SGA 4 (SGA 4 XVIII 1.4). In this letter to Larry Breen, I am assuming ``known'' the notion of \deleted{\unsure{a}}an $n$-stack (for $n=\deleted{\unsure{2|3}}\unsure{3|2}$ at any rate), and the corresponding notion of (strict) $n$-Picard stack, which should be describable (as was explained in Deligne's notes in the case $n=1$) by an $n$-truncated chain complex in the category of abelian sheaves on $\mathrm{X}$ (viewed mainly as an object of the relevant derived category). The ``strictness'' condition on usual\unsure{,|'} Picard stacks refers to the restriction that the commutativity isomorphism within an object $\mathrm{L}\otimes\mathrm{L}'$ , \deleted{\unsure{wh}} when $\mathrm{L}=\mathrm{L}'$, should reduce to the identity. It is assumed (wi thout further explanation) that the condition carries over in a natural way to Picard $n$-stacks, in such a way as to allow an interpretation of these by truncated objects in a suitable derived category, as hinted above.

(\textsuperscript{23}) When $\mathrm{M}$ is any abelian sheaf on a topos, $\mathrm{X}(\mathrm{M})$ is a certain canonical \inserted{left} resolution of $\mathrm{M}$ by sheaves of $\underline{\mathrm{Z}}$-modules which are ``free'', and more specifically, which are finite direct sums of sheaves of the type $\underline{\mathrm{Z}}^{(\mathrm{T})}$ , where $\mathrm{T}$ is any sheaf of the type $\mathrm{M}^{n}$ (finite product of copies of $\mathrm{M}$ ). This canonical\deleted{ly} construction was introduced \deleted{first} by MacLane, and gained new popularity in the French school of algebraic geometry and homological algebra in th\unsure{ê|e} late sixties, because \deleted{in case M is described by some abelian group scheme over a ground} it gives a very handy way to relate the $\mathrm{Ext}^{i}(\mathrm{M},\mathrm{N})$ invariants (when $\mathrm{N}$ is another abelian sheaf on $\mathrm{X}$) to the \deleted{\unsure{sp}} ``spacial'' cohomology of $\mathrm{M}$ (i.e. of the induced topos $\mathrm{X}_{/\mathrm{M}}$\deleted{\unsure{,}}) with coefficients in $\mathrm{N}$.

(\textsuperscript{24}) Reflecting on the ``right'' version of the provisional Verdier notion of a \deleted{derived}\inserted{triangulated} category (which was supposed to describe adequately the \inserted{relevant} internal structure of the derived categories of abelian categories) is part of my present program for the notes on Pursuing Stacks, and will be the main task in \deleted{\unsure{t}}one of the chapters of volume two.

(\textsuperscript{25}) As was seen in section 9, ``uniqueness'' here has to be understood in
% archive sheet number: 87
\scanpage{89}
\mspage{\handwritten{3\deleted{\unsure{4|7}}8}}
% A single diagonal crosses the typescript, including its footnote.
\deleted{au cours de cette année-là que j'ai appris pour la première fois (mais non \inserted{\handwritten{certes pour}} la dernière\handwritten{!}) \deleted{\unsure{certes|censés} \unsure{!|(*}}\deleted{)} que\deleted{\unsure{la}} c'est considéré comme un manque de sérieux de démontrer bel et bien des choses délicates et visiblement fondamentales, \inserted{\handwritten{que tout le monde utilise et}} que les prédécesseurs se sont toujours contentés d'admettre\inserted{\handwritten{(*)}}. Ou de démontrer un résultat qui englobe comme cas particuliers plusieurs résultats \inserted{profonds} connus (ce qui montre bien évidemment que le résultat soi disant nouveau ne peut être qu'un cas particulier ou une conséquence facil\unsure{e} de\deleted{\unsure{s}}s résultats connus). Ou de prendre la peine, \deleted{d'énoncer un}\inserted{dans \deleted{\unsure{un}}\handwritten{l'énoncé d'un}} résultat ou dans la description d'une situation en termes d'une autre, de formuler avec soin les hypothèses naturelles (signe d'un regrettable bombinage), plutôt que de se \deleted{contenter} borner \deleted{au cas simple} à quelque cas d'espèce au goût de celui qui émet son opinion. Ainsi j'ai vu faire reproche à Contou-Carrère de ne pas s'être borné dans sa thèse à se placer sur le corps des complexes ou tout au moins sur un corps, au lieu d'un schéma général\deleted{\unsure{.}} (tout en lui concédant \inserted{quand même} la circonstance atténuante que \deleted{\unsure{s}}\handwritten{c}'était \deleted{son patron \unsure{qui}} \inserted{\handwritten{sûrement sur les instances de son patron qu'il avait dû s'y résoudre)}}\handwritten{.} Celui qui s'exprimait ainsi était pourtant suffisamment dans le \deleted{\handwritten{\unsure{assister à quelque chose}}} coup pour savoir que même en se bornant au corps des complexes, les nécessités de\handwritten{s} \deleted{\unsure{la}} démonstration\handwritten{s} forçaient la main pour introduire des schémas de base générau\deleted{\unsure{n}}x ...\deleted{\unsure{?!|)}}}
\margin{\handwritten{\tikz[baseline]{\draw (0,0)--(.34,0)--(.40,.28)--(.73,.28); \draw (.31,0)--(.37,.28);}}}

\deleted{Les égarements d'une certaine mode \inserted{\handwritten{aujourd'hui}} vont jusqu'à \deleted{\unsure{mettre} exclure \unsure{tain|tair|tard|Les}}\inserted{honnir} non seulement \deleted{\unsure{d'une}}\inserted{\handwritten{\unsure{d|l}}es} démonstrations\handwritten{\ }en \deleted{bonne et due} forme, mais \inserted{\handwritten{souvent}} même \handwritten{d}es énoncés \deleted{en bonne ou en forme} et \handwritten{d}es définitions en forme\unsure{;|.|:} \deleted{\unsure{x}} au prix \deleted{d'un énoncé en bonne et due forme (sans même parler des hypothèses les plus naturelles), et} où est le papier et la longanimité du lecteur, il ne sera bientôt plus question de s'encombrer d'un luxe aussi coûteux\handwritten{!} \deleted{et} \deleted{on se bornera prochainement à}\inserted{\unsure{En} extrapolant les tendances actuelles}, il ne doit pas être difficile de \deleted{calculer}\inserted{\handwritten{prédire}} le moment exact où il ne sera plus question \inserted{dans une publication} d'expliciter \deleted{énoncés ni} définitions \inserted{\handwritten{ni énoncés}}, qu'on se contentera \deleted{de désigner}\inserted{désormais de nommer} par des mots-code, en laissant à l'infatigable et génial lecteur le soin de remplir les blancs, conformément à ses propres lumières. La tâ\deleted{\unsure{c}}che du referee sera facilitée d'autant, car il lui suffira de \deleted{consulter}\inserted{\handwritten{regarder}} dans l'annuaire \handwritten{``}Wo is Who\handwritten{''} si l'auteur est connu comme crédible (de toutes façons personne ne pourrait contredire les blancs et les \deleted{petits points}\inserted{\handwritten{pointillés}} qui composent le brillant article), ou au contraire un inavouable inconnu\handwritten{\ }qui sera (comme c'est le cas déjà dès aujourd'hui) \deleted{éliminé}\inserted{\handwritten{éjecté}} d'office ...}

\noindent\rule{2em}{.4pt}\par
% slanted
\margin{\handwritten{\unsure{sans interligne|sans interlignes}}}
\deleted{(\textsuperscript{*}) Voir aussi à ce sujet l'épisode ``La note - ou la nouvelle éthique'' \handwritten{(}par\unsure{.}33 \handwritten{)}. \deleted{Qui}\inserted{\handwritten{Cette fameuse ``note''}} avait justement le tort d'expliciter des notions \deleted{ou}\inserted{\handwritten{et}} des \deleted{dém} énoncés \deleted{et des démonstrations} qui avaient été laissées jusque là dans le vague, et qui pourtant ont été \deleted{utilisés (par moi en \unsure{tout autant|tout allant} par \unsure{les|le} \unsure{élèves|élevé|évidé}) sans vergogne}\inserted{implicitement utilisés par moi} pour établir des résultats qui portent mon nom et \deleted{\unsure{que}} que tout le monde \deleted{(y compris \inserted{\handwritten{\unsure{mon}}}\unsure{maître et élève})} utilise sans vergogne\deleted{\unsure{.|,}}\inserted{\handwritten{depuis bientôt vingt-cinq ans.}}}
% archive sheet number: 88
\scanpage{90}
\mspage{\handwritten{65}}

considerably wider sense than I expected, when writing this letter to Larry Breen. It now appears that the whole theory of stacks of groupoids will depend on the choice of a ``coherator'' $\underline{\mathrm{C}}$ , as seen in section 13.

(\textsuperscript{26}) Tim Porter pointed out to me that ``Dold Puppe'' is an accurate name for this basic theorem, which should be called Dold-Kan theorem.

(\textsuperscript{27}) As was pointe\deleted{\unsure{e}}d out to me by Ronnie Brown, this structure was already well-known to J.H.C. Whitehead, under the name of ``crossed modules'', and extensive use and extensive generalizations of this notion (in quite different directions from those I was having in mind, in terms of Gr-stacks over an arbitrary topos) have been made by him and others. With respect to the question on next page, of generalizing this notion of ``non commuative chain complex'' from length one to length two, Ronnie says there is a work in preparation by D. Conduché `` Modules croisés généralisés de longueur 2 '' .

(\textsuperscript{28}) Tim P\deleted{\unsure{e}}orter pointed out to me that work on étale tubular neighbourhoods was done by D.A. Cox ; ``Algebraic tubular neighbourhoods I,II', Math. Scand. 42 (1978) 211-228, 229-242. I've not seen yet this work, and can't say therefore whether it meets the rather precise expectations I have for a theory of tubular neighbourhoods, for the needs of a dévissage theory of stratified schemes (or, more generally, stratified topoi).

(\textsuperscript{29}) This is the typical game \deleted{to be} embodied in the ``derivator'' associa\deleted{\unsure{te}}ted to the theory $(\mathrm{Hot})$ of usual homotopy types (compare section 69).

(\textsuperscript{30}) Some more details on this program are outlined in Esquisse d'un \deleted{\unsure{P}}programme (section 5), in Réflexions Mathématiques 1 .

(\textsuperscript{31}) I was informed by knowledgeable people soon later that the answer is well known to be negative, by working with ``rational homotopy types'' (the cohomology of which is made up with vector spaces over $\underline{\mathrm{Q}}$ . It is well known indeed that a 1-connected rational homotopy type is not known from its rational cohomology ring alone, which is supposed (I confess I did'nt check) to contain already all the information I was contemplating.

(\textsuperscript{32}) This \deleted{is the \unsure{``now standard''|now standard} version of the similar problem}\inserted{``problem'' is met with by the notion of a ``derivator''} , which ``was in the air'' already by the late sixties, but was never developped - (instead even derived categories became tabu in the seventies ...). Compare with note (\textsuperscript{24}).
% archive sheet number: 89
\scanpage{91}
% Printed seminar programme, rotated 180 degrees relative to the scan.
% A diagonal stroke rules out the entire printed page.
\deleted{SÉMINAIRE NICOLAS BOURBAKI}

\deleted{RÉSUMÉS DES EXPOSÉS}

\deleted{BEAUVILLE Arnaud \quad \textit{Surfaces} K3}

\deleted{La démonstration récente par Y.T. Siu du fait que toute surface K3 est Kählérienne achève une série de travaux qui ont permis de décrire complètement ``l'espace des modules'' des surfaces K3 en termes de leurs périodes. On se propose de décrire les grandes lignes de cette théorie : injectivité de l'application des périodes, surjectivité, cône des classes de Kähler.}

\deleted{CHAPERON Marc \quad \textit{Quelques questions de géométrie symplectique [d'après, entre autres, Poincaré, Arnold, Conley et Zehnder]}}

\deleted{Le dernier théorème géométrique de Poincaré - démontré par G.D. Birkhoff - a conduit V.I. Arnold à formuler de très intéressants problèmes globaux en géométrie symplectique. Dans un travail récent, C.C. Conley et E. Zehnder utilisent une généralisation naturelle de la théorie de Morse pour résoudre l'un d'entre eux, concernant le nombre de points fixes des transformations symplectiques ``homologues à l'identité'' d'une variété symplectique compacte. Une modification presque immédiate de leur méthode permet de répondre à une question voisine sur les intersections de variétés lagrangiennes.}

\deleted{DOLGACHEV Igor \quad \textit{Integral quadratic forms : Algebro-geometrical applications [after V. Nikulin]}}

\deleted{In this talk we will report on the work of V. Nikulin on arithmetics of integral quadratic non-unimodular forms and their applications to surface singularities, automorphisms of K3-surfaces and topology of real algebraic surfaces.}

\deleted{DUFLO Michel \quad \textit{Analyse harmonique sur les groupes algébriques complexes : Formule de Plancherel [d'après M. Andler] et conjecture de M. Vergne}}

\deleted{La méthode des orbites, inaugurée par A.A. Kirillov dans son étude des groupes de Lie nilpotents, permet de décrire un certain nombre des objets traditionnels de l'analyse harmonique sur un groupe de Lie en termes de la représentation co-adjointe. Pour les groupes algébriques complexes, on connait ainsi la formule de Plancherel, et on peut résoudre une conjecture de M. Vergne sur la formule de ``Poisson-Plancherel''.}

\deleted{HAKIM Monique \quad \textit{Valeurs au bord de fonctions holomorphes bornées en plusieurs variables complexes}}

\deleted{Il existe des fonctions holomorphes bornées non constantes dans la boule unité de $\mathbb{C}^{n}$, dont les limites radiales sont presque partout de module 1. On expose ce résultat dû à A.B. Aleksandrov et on étudie des constructions de fonctions holomorphes bornées ayant des propriétés remarquables, introduites par divers auteurs.}

\deleted{SADOURNY Robert \quad \textit{Quelques aspects mathématiques de la dynamique des fluides géophysiques}}

\deleted{La dynamique des fluides géophysiques se partage entre l'expérimentation numérique lourde et quelques approches plus théoriques : existence d'attracteurs dans les espaces de solutions, permettant l'approche de certains mécanismes au moyen d'équations filtrées ; étude, par les méthodes des systèmes dynamiques, des changements de régime de la circulation atmosphérique ou du climat.}
% archive sheet number: 90
\scanpage{93}
\mspage{\handwritten{1}}
% Only the bottoms of strokes survive at the top right of this copy.
\margin{\illegible{3}}

\noindent Dear Daniel,\\
\margin{\handwritten{1.}}Last year Ronnie Brown from Bangor sent me a heap of reprints and preprints by him and a group of friends\handwritten{,} on various foundational matters of homotopical algebra. I did not really dig through any of this, as I kind of lost contact with the technicalities of this kind (I was never too \unsure{familiar} with the homotopy techniques \inserted{\handwritten{anyhow,}} I confess) - but this reminded me of a few letters I had exchanged with Larry Breen in\unsure{.}1975, where I had developped an outline of a \inserted{program for a} kind of ``topological algebra'', viewed as a synthesis of \deleted{\unsure{ho}}\unsure{homotopical} and homological algebra, with special emphasis on topoi - most of the basic intuitions in this program arising from various \unsure{backgrounds|background} in algebraic geometry. \deleted{(}Some of those intuitions we discussed\handwritten{,} I believe\handwritten{,} at IHES eight or nine years before, at a time when you had just written up your nice ideas on axiomatic homotopical algebra, published since in Spinger's Lec\unsure{t}ure Notes. I write you under the assumption that you have not entirely lost interest for those foundational questions you were looking at more than fifteen years ago. One thing which strikes me, is that (as far as I know) there has not been any substantial progress since - it l\deleted{\unsure{a}}ooks to me that \deleted{\unsure{adequate}} understanding of the basic structures underlying homotopy theory, or even homological algebra only, is still \unsure{lacking} - \unsure{probably} because the few people who have a wide enough background and perspective enabling them to feel the main questions, are devoting their energies to things which seem more directly rewarding. Maybe \deleted{there} even a wind of disrepute for any foundational matters whatever is blowing nowadays ! In this respect, what seems to me even more striking than the lack of proper foundations \handwritten{f}or homological and homotopical algebra, is the absence I dare\unsure{-}say\unsure{,} of proper foundations for topology itself ! I am thinking here \unsure{mainly} of the development of a context of ``tame'' topology, which (I am \unsure{convinced)} would have on the everyday technique of geometric topology (\deleted{in} I use this expression in contrast to the topology of use for analysts) a comparable impact or even a greater one, than the introduction of the point of view of schemes had on algebraic geometry. The \inserted{psychological} drawback here I believe is not anything like messyness, as for homological and homotopical algebra\inserted{\handwritten{\unsure{(or for schemes)}}}, but \unsure{merely} the inrooted intertia which \deleted{\unsure{goes against}} prevents \deleted{\unsure{everyone}} us so stubbornly from looking innocently, with fresh eyes, upon things, without being dulled and emprisoned by standing habits of thought, going with a familiar context - \uline{too} familiar a context ! The task of working out the foundations of tame topology\handwritten{,} and a \deleted{theory} corresponding structure theory for ``stratified (tame) spaces\handwritten{''}, seems to me a lot more urgent and exciting still \unsure{than}
% archive sheet number: 92
\scanpage{94}
\mspage{\handwritten{1'}}

any program of homological, homotopic\handwritten{\unsure{a}}l or topological algebra.

\margin{\handwritten{2.}}The motivation for this letter was the latter topic however. \deleted{\unsure{Br}}\unsure{Ronnie} Brown and his friends are competent algebraists and apparently strongly motivated for investing energy in foundational work, on the other hand they visibly are lacking the necessary scope of vision which geometry alone provides. They seem to me kind of\unsure{:}isolated, partly due I guess to the disrepute I mentioned before \unsure{o-|-} I suggested to try and have contact with people such as yourself, Larry Breen, Illusie and others, who have the geometric insight and \deleted{\unsure{what}}w\handwritten{ho} moreover, \deleted{I believe} may not think themselves to ggod for indulging in occasional reflection o\handwritten{n} foundational matters\handwritten{,} and in the process help others do the work which should be done.

At first sight it had \deleted{\unsure{appea}}\unsure{seemed|appead} to me that the Bangor group had indeed come to work out (quite independently) one basic intuition of the program I had envisioned in those letters to Larry Breen - namely that the study of $n$-truncated homotopy types (of semisimplicial sets, or of topological spaces) was essentially equivalent to the study of so-called $n$-groupoids (where $n$ is any natural integer). This is \unsure{ex$\hat{\mathrm{p}}$ected|expected} to be \margin{\handwritten{\unsure{``}}}achieved by associating to any space (say) $\mathrm{X}$ it's``fundamental $n$-groupoid'' $\handwritten{\mathbb{\mathord{\rlap{\kern.23em\vrule height .66em width .035em}\mathrm{\Pi}}}}_{n}(\mathrm{X})$, generalizing the familiar Poincaré fundamental groupoid for $n=1$. \margin{\handwritten{\unsure{'}}}The o\unsure{B|b}vious idea is that 0-objects of $\handwritten{\mathbb{\mathord{\rlap{\kern.23em\vrule height .66em width .035em}\mathrm{\Pi}}}}_{n}(\mathrm{X})$ should be the points of $\mathrm{X}$, 1-objects should be ``homotopies'' or paths between points, 2-objects should be homotopies between 1-objects, etc.\unsure{-} This $\handwritten{\mathbb{\mathord{\rlap{\kern.23em\vrule height .66em width .035em}\mathrm{\Pi}}}}_{n}(\mathrm{X})$ should embody the $n$-truncated homotopy type of $\mathrm{X}$, in much the same way as for $n=1$ the usual fundamental groupoid embodies the 1-truncated homotopy type. For two spaces $\mathrm{X},\mathrm{Y}$, the set of homotopy-classes of maps $\mathrm{X}\handwritten{\to}\mathrm{Y}$ (more correctly, for general $\mathrm{X},\mathrm{Y}$, the maps of $\mathrm{X}$ into $\mathrm{Y}$ in the homotopy category) should correspond to $n$-equivalence classes of $n$-functors from $\handwritten{\mathbb{\mathord{\rlap{\kern.23em\vrule height .66em width .035em}\mathrm{\Pi}}}}_{n}(\mathrm{X})$ to $\handwritten{\unsure{\mathbb{\mathord{\rlap{\kern.23em\vrule height .66em width .035em}\mathrm{\Pi}}}}}_{n}(\mathrm{Y})$ - etc. There are very stro\unsure{-}ng suggestions for a nice formalism, including a notion of geometric realization of an $n$-groupoid, which should imply that any $n$-groupoid is $n$-equivalent to a $\handwritten{\mathbb{\mathord{\rlap{\kern.23em\vrule height .66em width .035em}\mathrm{\Pi}}}}_{n}(\mathrm{X})$. Moreover, when the notion of $n$-groupoid (or more generally of an $n$-category) is relativized over an arbitrary topos to the notion of an $n$-gerb (or more generally, an $n$-stack), these become the \deleted{\unsure{main}} natural ``coefficients'' for a formalism of non commutative cohomol\unsure{o}gical algebra, in the spirit of Giraud's thesis.

But all this kind of thing\unsure{s} for the time being is pure heuristics - I never so far sat down to try to make explicit at least a definition of $n$-categories and $n$-groupoids, of $n$-functors between these etc. When I got the Bangor reprints I at once had the feeling that this kind of \handwritten{work}
% The extra upright in each pen Pi is drawn explicitly for the supplied math fonts.
% At Pi_n(Y) the top stroke is lost in this copy.
% archive sheet number: 93
